\documentclass[10pt,a4paper]{cv}
\geometry{left=1cm,right=11.2cm,marginparwidth=9cm,marginparsep=1.2cm,top=1cm,bottom=1cm}

%\usepackage[T2A, T1]{fontenc}


\usepackage[utf8]{inputenc}
\usepackage[russian]{babel}
\usepackage[default]{lato}
\usepackage{hyperref}
\fontfamily{sans-serif}

\definecolor{VividPurple}{HTML}{0085c3}
\definecolor{SlateGrey}{HTML}{2E2E2E}
\definecolor{LightGrey}{HTML}{666666}
\colorlet{heading}{VividPurple}
\colorlet{accent}{VividPurple}
\colorlet{emphasis}{SlateGrey}
\colorlet{body}{LightGrey}

\usepackage{xcolor}


\hypersetup{
    colorlinks,
    linkcolor={VividPurple},
    citecolor={VividPurple},
    urlcolor={VividPurple}
}


% itemize and rating marker for \cvskill 
\renewcommand{\itemmarker}{{\small\textbullet}}
\renewcommand{\ratingmarker}{\faCircle}


%\usepackage{hyperref}

\begin{document}
%
\name{Александр Галайда}
\tagline{Резюме}
\photo{3.5cm}{face}
%
\personalinfo{%
  \begin{minipage}{4cm}
  	\email{\href{mailto:alexgalayda@gmail.com}{\textcolor{SlateGrey}{alexgalayda@gmail.com}}}
   
   \github{\href{https://github.com/alexgalayda}{\textcolor{SlateGrey}{github.com/alexgalayda}}}
   
   \github{\href{https://gitlab.com/alexgalayda}{\textcolor{SlateGrey}{gitlab.com/alexgalayda}}}
   
   \end{minipage}%
   \hfill%
   \begin{minipage}{5.5cm}
   
   \phone{+7-916-395-20-46}
   \location{Москва, Россия}
  \end{minipage}
}
	

\begin{adjustwidth}{}{-10cm}
\makecvheader
\end{adjustwidth}

\cvsection[page1sidebar_ru]{ОПЫТ РАБОТЫ}

% Физтех
\cvevent{\textbf{Разработчик-исследователь алгоритмов}}{\href{https://mipt.ru/}{МФТИ (ФРТК), Кафедра Радиоэлектроники и Прикладной Информатики}}{11.18 -- по н. в.}{Долгопрудный, Россия}
Разработка ПО убирающего артефакты с изображений, ПО для геопривязки изображений, ПО с использованием нейросетей для детектирования аномалий на снимках с беспилотных утройств

\divider

% Аналитик
\cvevent{\textbf{Рисковый аналитик}}{\href{www.tinkoff.ru}{Тинькофф Банк}}{08.18 -- 11.18}{Москва, Россия}
Анализ рисков, построение скоринговой модели страховок

\divider

% Учитель
\cvevent{\textbf{Учитель}}{\href{https://sch1298sz.mskobr.ru/\#/}{Школа №1298 Куркино}}{02.18 -- 06.18}{Химки, Россия}
Вел курсы по олимпиадной математике для старшеклассников

\divider

\cvsection{ДОСТИЖЕНИЯ}
\cvsubsection{Международный форум \href{https://www.rusarmyexpo.ru/}{<<Армия 2019>>}}
Доклад на тему <<Обзор способов формирования и различных приемов расширения обучающей выборки для
искусственных нейронных сетей>>

\cvsubsection{Стипендия им. Абрамова за успехи в учёбе в университете в 2015 - 2016г.}


\cvsection{ПРОЙДЕННЫЕ КУРСЫ}
\cvsubsection{Операционная система Линукс и Линукс-серверы }{
М.Н.Штейнберг, заведующий лабораторией компьютерных и интернет технологий, \href{http://geo.mipt.ru/}{СУМГФ}, МФТИ}

\cvsubsection{Машинное обучение и анализ данных}{
	Специализация на \href{https://ru.coursera.org/specializations/machine-learning-data-analysis}{Coursera от МФТИ и Яндекса}}
%\cvsubsection{Kalman and Bayesian Filters in Python}{курс \href{https://github.com/rlabbe/Kalman-and-Bayesian-Filters-in-Python}{Roger R. Labbe}}

%\cvsubsection{Построение глубоких нейронных сетей на языке Python}{
%	курс \href{https://www.asozykin.ru/courses/nnpython}{Андрея Созыкина}}


\clearpage
\end{document}
